\documentstyle[% 


righttag% 


,11pt% 


,amssymb% 


,russian%               remove in Eng. 


,izvru%                 Set izven% in Eng. 


% ,izvru-ks             for brief communications 


]{amsart} 


 


\theoremstyle{definition} 


\theorembodyfont{\rm} 


\newtheorem{defnnonum}{\hskip\parindent Определение} 


\renewcommand{\thedefnnonum}{} 


 


%\hoffset=-15mm                  For Eng. 


 


\setcounter{page}{77} 


\god{1997} 


\nomer{8~(423)} %                Remove in Eng. 


%\nomer{Vol.\,41, No.\,7}        Set in Eng. 


% \str{75-81}                    Set in Eng. 


\udk{517.988} 


\author{\it Н.М.\,Ратинер} 


\title{Свойства классов бордизмов оснащенных\\


подмногообразий фредгольмовой структуры} 


 


\thanks{Работа выполнена при поддержке Российского фонда фундаментальных


исследований, грант \No\,94-01-00183.} 


\date{} 


% \pagestyle{myheadings}         Set in Eng. 


 


\pagestyle{plain} %              Remove in Eng. 


\begin{document} 


 


% \thispagestyle{plain}          Set in Eng. 


 


\maketitle 


 


% Body text 


 


Оснащенные подмногообразия фредгольмовой структуры на банаховом 


многообразии $X$ и группы их бордизмов появились в работе Элворти и 


Тромба \cite{Elw} при определении степени фредгольмовых отображений 


положительного индекса, а затем были использованы в работах В.Г.\,Звягина и 


автора \cite{Zv1}, \cite{Zv2} для обобщения степени на более широкие 


классы отображений 


и приложений теории степени к исследованию задач разрешимости и 


бифуркации решений нелинейных дифференциальных уравнений. 


 


Данная статья посвящена изучению свойств классов бордизмов оснащенных 


подмногообразий фредгольмовой структуры $X_\Phi$ на банаховом многообразии. 


Проблемы, исследуемые в настоящей работе, возникли при использовании 


теории степени фредгольмовых отображений положительного 


индекса в \cite{Zv1}, \cite{Zv2} и нашли частичное решение в 


\cite{Zv2}. 


 


{\bf 1.} Введем некоторые обозначения и определения. Пусть $E$ --- 


вещественное 


банахово пространство, $GL_c(E)$ --- множество линейных обратимых 


операторов вида $I+K$, где $I$ --- тождественный, а $K$ --- вполне 


непрерывный 


операторы. 


 


Пусть $X$ --- банахово многообразие. Фредгольмовой структурой $X_\Phi$ 


называется максимальный атлас $X_\Phi=\{(U_i,\varphi_i)\}_i$ со свойством: 


для любых 


двух карт $(U_i,\varphi_i)$ и $(U_j,\varphi_j)$, 


$U_i\cap U_j\ne\emptyset$ и $x\in U_i\cap U_j$, 


$d(\varphi_j\cdot\varphi_i)\varphi_i(x)\in GL_c(E)$.


 


Пусть $\eta:E\to M$ --- векторное расслоение со слоем $E$ над 


многообразием 


$M$ (банаховым или конечномерным). Расслоение $\pi$ называется 


$GL_c$-расслоением, 


если его структурная группа редуцирована к $GL_c(E)$, т.\,е. для 


тривиализующих 


отображений $\{\tau_i:\pi^{-1}(U_i)\to U_i\times E\}_i$ отображения 


перехода 


$\tau_i\cdot\tau_j^{-1}:U_i\times E\to U_j\times E$ имеют вид 


\begin{equation} 


(x,v)\to(x,v+a(x)v), 


\end{equation} 


где $a(x)$ --- вполне непрерывный линейный оператор при каждом $x$. 


 Если в 


формуле (1) $a(x)$ --- локально конечномерный оператор, то расслоение 


$\pi$ 


называется $l$-расслоением. 


 


Изоморфизм $\xi:\pi_1\to\pi_2\;$ $GL_c$-расслоений называется 


$GL_c$-изоморфизмом, 


если в локальной записи он имеет вид $(x,v)\to(x,v+b(x)v)$, где 


$b(x)$ --- вполне непрерывный линейный оператор при каждом $x$; $\xi$ 


называется 


$l$-изоморфизмом, если $b(x)$ локально конечномерный. 


 


Будем считать в дальнейшем, что модельное пространство фредгольмовой 


структуры $X_\Phi$ есть $E\times R^q$. 


 


Пара 


$(M^q,i)$, где $M^q$ --- замкнутое $q$-мерное гладкое многообразие, а 


$i:M^q\to X$ --- непрерывное отображение, будет называться сингулярным 


подмногообразием в $X$. $GL_c$-оснащением пары $(M,i)$ в $X_\Phi$ будет 


называться 


$GL_c$-изоморфизм векторных расслоений над 


$M:\tau:TM\oplus(M\times E)\to i^*(TX_\Phi)$, где 


$TM$ --- касательное расслоение к $M$, $i^*(TX_\Phi)$ --- обратный образ 


касательного 


расслоения к структуре $X_\Phi$. Обозначим через $S_q(X_\Phi)$ множество 


троек 


$(M,i,\tau)$, где $(M,i)$ --- $q$-мерное сингулярное многообразие в $X$, 


а $\tau$ 


--- его 


оснащение. Введем в $S_q(X_\Phi)$ следующее отношение эквивалентности, 


называемое 


отношением $GL_c$-оснащенной бордантности. 


 


\begin{defnnonum} Тройки $(M_0,i_0,\tau_0)$ и $(M_1,i_1,\tau_1)$ 


называются $GL_c$-оснащенно бордантными, 


если существует тройка $(Y,i,\tau)$, где $Y$ --- $(q+1)$-мерное 


компактное многообразие с краем $\partial Y=M_0\cup M_1$ ($Y$ --- 


``пленка'', соединяющая 


$M_0$ и $M_1$), $i:Y\to X\times[0,1]$ --- непрерывное отображение такое, 


что 


$\left. i\right|_{M_r}=i_r$ ($r=0,1$), и $\tau$ --- $GL_c$-изоморфизм 


расслоений: $\tau:TY\oplus(Y\times E)\to i^*T(X_\Phi\times[0,1])$, 


совпадающий с $\tau_r$ на $TM_r\oplus(M_r\times E)$, $r=0,1$. 


\end{defnnonum} 


 


Фактор-множество $S_q(X_\Phi)$ по отношению $GL_c$-оснащенной бордантности 


обозначается $F_q(X_\Phi)$ и называется группой $GL_c$-оснащенных бордизмов. 


Класс эквивалентности тройки $(M,i,\tau)$ в $F_q(x_\Phi)$ обозначается 


через 


$[M,i,\tau]$. 


 


Напомним, что нормальным расслоением иммерсии $i:M\to\tau$ называется 


фактор-расслоение $i^*(TX_\Phi)/TM=\nu_X M$. Предположим, что расслоение 


$\nu_X M$ тривиализуемо, 


т.\,е. имеется $GL_c$-изоморфизм $\xi:M\times E\to\nu_X M$. Тогда пара 


$(M,i)$ 


имеет $GL_c$-оснащение вида 


\begin{equation} 


\tau=I_{TM}\oplus\xi:TM\oplus(M\times E)\to TM\oplus\nu_XM\approx 


i^*(TX_\Phi), 


\end{equation} 


которое будем называть правильным оснащением. 


 


Оснащения вида (2) представляют собой обобщение конструкции 


Л.С.\,Понтрягина \cite{Pon}, который оснащением подмногообразия $M^q$ в 


конечномерном 


евклидовом пространстве $R^{n+q}$ называл базис в нормальном пространстве 


к $M^q$, непрерывно зависящий от точки $x\in M^q$. 


 


{\bf 2.} Сформулируем основные теоремы. 


 


\begin{thm} {\em (а)} Если банахово многообразие $X\;$ $q$-связно, то для 


любого 


элемента $\alpha$ группы $F_q(x_\Phi)$ найдется его представитель 


$(M,i,\tau)$ такой, 


что $i:M\to U\subset X$ --- вложение, где $(U,\varphi)$ --- некоторая 


карта структуры 


$X_\Phi$, причем образ $\varphi\cdot i(M)$ конечномерен; 


 


{\em (б)} если банахово многообразие $X\;$ $(q+1)$-связно, 


$(M_0,i_0,\tau_0)$, 


$(M_1,i_1,\tau_1)$ --- две $GL_c$-оснащенно бордантные тройки такие, как 


в {\em (а)}, то 


пленка $(Y,i,\tau)$, соединяющая их, может быть выбрана так, что 


$i:Y\to U\times[0,1]\subset X\times[0,1]$ ---


вложение, а образ $(\varphi\times I)\cdot i(Y)$ конечномерен. 


\end{thm} 


 


Говорят, что банахово пространство $E$ обладает свойством конечномерной 


аппроксимации, если для каждого банахова пространства $F$ множество 


линейных непрерывных конечномерных операторов из $E$ в $F$ плотно в 


множестве линейных вполне непрерывных операторов. В частности, 


пространство с базисом Шаудера обладает свойством конечномерной 


аппроксимации (\cite{Lin}). 


 


Следующая теорема показывает, что каждое $GL_c$-оснащение можно заменить 


на надстройку над конечномерным оснащением, не выходя из класса 


$GL_c$-оснащенных бордизмов. 


 


\begin{thm} Пусть $X\subset E\times R^q$ --- область в банаховом 


пространстве, 


обладающем свойством конечномерной аппроксимации. Тогда для каждого 


$\alpha$ из 


$F_q(X_\Phi)$ найдется тройка $(M,i,\tau)\in\alpha$, где оснащение $\tau$ 


имеет вид 


\begin{multline} 


TM\oplus(M\times E)=TM\oplus(M\times E^n)\oplus 


(M\times E^{\infty-n})@>{\tau_n\oplus I_{M\times E^{\infty-m}}}>> \\


%


(i(M)\times E^n\times R^q)\oplus(i(M)\times E^{\infty-n}), 


\end{multline} 


$E^n$ и $E^{\infty-n}$ --- два дополнительных замкнутых подпространства в 


$E$, 


$\dim E^n<\infty$, а 


$\tau_n:TM\oplus(M\times E^n)\to(i(M)\times E^n\times R^q)$ --- 


конечномерное оснащение. 


Пленку $(Y,i,\tau)$, соединяющую две тройки с оснащениями вида {\em (3)}, 


можно 


выбрать так, чтобы $GL_c$-оснащение $\tau$ имело вид 


$$ 


\tau:TY\oplus(Y\times E)=TY\oplus(Y\times E^m)\oplus 


(Y\times E^{\infty-m}) 


@>{\tau_m\oplus I_{Y\times E^{\infty-n}}}>>i(Y)\times E^m 


\times R^q\times E^{\infty-m}, 


$$ 


где $E^m$, $E^{\infty-m}$ --- два дополнительных замкнутых подпространства 


в $E$, $\dim E^m<\infty$, $E^n\subset E^m$, 


а $\tau_m$ --- конечномерное оснащение. 


\end{thm} 


 


\begin{thm} {\em (а)} Если банахово многообразие $X\;$ $q$-связно, 


а его модельное пространство обладает свойством конечномерной 


аппроксимации, 


то для любого класса $\alpha$ из $F_q(X_\Phi)$ найдется такой его 


представитель 


$(M,i,\tau)$, что $i$ --- вложение, а $\tau$ --- правильное оснащение. 


 


{\em (б)} если $Z$ --- $(q+1)$-связное банахово многообразие, модельное 


пространство 


которого обладает свойством конечномерной аппроксимации, 


а тройки $(M_0,i_0,\tau_0)$ и $(M_1,i_1,\tau_1)$ такие, 


как в {\em (а)}, то


пленка $(Y,i,\tau)$, 


соединяющая эти две тройки, может быть выбрана так, что $i$ 


--- вложение, а 


$\tau$ --- правильное оснащение. 


\end{thm}


 


{\bf 3.} Изложим кратко схемы доказательства теорем. 


Для доказательства теоремы 1 понадобятся две леммы. 


 


\begin{lem} Если $(M,i)$ --- сингулярное подмногообразие в $X_\Phi$ с 


$GL_c$-оснащением, 


и отображение $i_1:M\to X$ непрерывно гомотопно $i$, то для пары 


$(M,i_1)$ существует $GL_c$-оснащение $\tau_1$ и тройки 


$(M,i,\tau)$ и $(M_1,i_1,\tau_1)$ 


$GL_c$-оснащенно бордантны. 


\end{lem} 


 


\begin{pf*}{Доказательство леммы 1} Для гомотопных отображений $i$ и 


$i_1$, 


индуцированные 


расслоения $i^*(TX_\Phi)$ и $i^*_1(TX_\Phi)\;$ $GL_c$-изоморфны, что дает 


оснащение 


для $(M,i_1)$. Для построения пленки, соединяющей $(M,i)$ и 


$(M,i_1)$ используется 


теорема 


о накрывающей гомотопии для расслоений со структурной группой 


$GL_c(E)$ (см. \cite{Hm}, с.\,79). 


\end{pf*} 


 


\begin{lem} Если $X$ есть $q$-связное банахово многообразие, а $M$ есть 


$q$-мерное компактное многообразие, то любое непрерывное отображение 


$i:M\to X$ гомотопно постоянному. 


\end{lem} 


 


Для доказательства леммы 2 многообразие $M$ рассматривается как 


$CW$-комплекс. 


Гомотопия задается последовательно по $k$-мерным остовам, где 


$0\le k\le q$. $\quad\square$ 


 


\begin{pf*}{Доказательство теоремы 1} Учитывая леммы 1 и 2, можем считать, 


что образ $i(M)$ лежит в одной карте $(U,\varphi)$ структуры $X_\Phi$, где 


$U$ --- стягиваемое 


подмножество в $X$. Для простоты, заменяя $i$ на $\varphi\cdot i$, считаем 


$X$ открытым 


подмножеством в модельном пространстве $E\times R^q$. Заметим, что 


$i(M)$ --- компактное подмножество в $X\subset E\times R^q$, 


следовательно, существует 


конечномерное подпространство $E^m\subset E\times R^q$ и непрерывный 


проектор $p:E\times R^q\to E^m$ 


такой, что $\|p\cdot i(x)-i(x)\|<\varepsilon$ для $x\in M$ (см. \cite{Kra}, 


с.\,123). Осталось воспользоваться 


тем, что множество вложений плотно в $C^1(M,E^m)$ и, следовательно, 


в любой $\sigma$-окрестности отображения $p\cdot i$ найдется вложение 


$j:M\to E^m$. 


Если $\varepsilon$ и $\delta$ достаточно малы, то $j(M)\subset X$. 


\end{pf*} 


 


Для доказательства теоремы 2 понадобится следующая лемма, доказательство 


которой использует свойство конечномерной аппроксимации для 


слоя расслоений. 


 


\begin{lem} Пусть $h$ есть $GL_c$-изоморфизм $l$-расслоений над 


конечномерным 


компактным многообразием $Y$. Тогда найдется $l$-изоморфизм $h_1$ такой, 


что $h_2=\mu_1(x)h+\mu_2(x)h_1$ является $l$-изоморфизмом для любых 


неотрицательных функций 


$\mu_1(x)$ и $\mu_2(x)$ на $Y$ с условием


$\mu_1(x)+\mu_2(x)\equiv 1$. Если $Y$ -- 


многообразие с краем и $\left. h\right|_{\partial Y}$ 


есть $l$-изоморфизм, то 


$\left. h_2\right|_{\partial Y}=\left. h\right|_{\partial Y}$. 


\end{lem} 


 


\begin{pf*}{Доказательство леммы 3} В каждом элементе $V_i$ 


тривиализующего покрытия 


выбирается точка $x_i$ и вполне непрерывный оператор $a(x_i)$ 


аппроксимируется 


конечномерным $a_i(x_i)$. Локальные изоморфизмы расслоений 


$(x,v)\to(x,v+a_i(x_i)v)$ склеиваются с помощью разбиения единицы на 


компактном 


конечномерном многообразии $Y$. 


\end{pf*} 


 


\begin{pf*}{Доказательство теоремы 2} Расслоения, в которых действует 


$GL_c$-изоморфизм, задающий оснащение, являются $l$-расслоениями (ввиду 


того, 


что $X$ --- область в $E\times R^q$). Из леммы 3 вытекает, что найдется 


покрытие 


$\{V_j\}$ базы расслоения такое, что оснащение в локальной записи имеет 


вид $(x,v)\to(x,v+a_j(x)v)$, где 


$\text{Im}\,a_j(x)\subset E_j^m\times R^q$, $\dim E_j^m<\infty$. Это 


покрытие можно считать конечным в силу 


компактности базы. Тогда 


$\tau_m=\left.\tau\right|_{TM\oplus(M\times E^m)}$ --- изоморфизм 


конечномерных векторных расслоений и 


$\tau=\tau_m\oplus I_{M\times E^{\infty-m}}$, где 


$E_j^m\subset E^m$ для всех $j$. 


\end{pf*} 


 


{\bf Доказательство теоремы 3} опирается на следующую лемму о приведении 


оснащения $q$-мерного подмногообразия $(n+q)$-мерного евклидова 


пространства 


к правильному виду, но в $2(n+q)$-мерном пространстве. 


 


\begin{lem} Пусть $M$ --- $q$-мерное подмногообразие в евклидовом 


пространстве 


$E^{n+q}=E^n\times E^q$ с оснащением 


$\tau:TM\oplus(M\times E^n)\to M\times E^{n+q}$. Тогда оснащение 


$\tau\oplus I_{M\times E^{n+q}}:TM\oplus(M\times E^{n+q}) 


\oplus(M\times E^{n+q})\to M\times E^{2(n+q)}$ 


эквивалентно правильному оснащению 


$I\oplus\xi$, где $\xi:M\times E^{2(n+q)}\to\nu_{2(n+q)}M$ --- 


тривиализация нормального 


расслоения к $M$ в $E^{2(n+q)}$. 


\end{lem} 


 


\begin{pf*}{Доказательство леммы 4} Заменим сначала оснащение $\tau$ 


эквивалентным 


ему оснащением, для которого образы $\tau(TM)$ и $\tau(M\times E^n)$ 


ортогональны. Для 


изоморфных расслоений $TM$, $\tau(TM)$ классифицирующие отображения из 


многообразия 


$M$ в грассманово многообразие $G_q^{2(n+q)}$ гомотопны 


(\cite{Hm}, с.\,53). 


Обратный образ 


канонического расслоения $\gamma_q^{2(n+q)}$ над $G_q^{2(n+q)}$ при этой 


гомотопии можно считать подрасслоением $\eta$ в 


$(M\times[0,1])\times E^{2(n+q)}$, которое 


совпадает с $TM$ над $M\times\{0\}$ и с $\tau(TM)$ над 


$M\times\{1\}$. Его ортогональное 


дополнение $\eta^{\perp}$ 


совпадает с $\nu_{2(n+q)}M$ над $M\times\{0\}$ и с 


$\tau(M\times E^n)\oplus(M\times E^{n+q})$ над 


$M\times\{1\}$. Для расслоений $\eta$ и $\eta^{\perp}$ над цилиндром 


$M\times[0,1]$ 


имеются изоморфизмы 


$g:TM\times[0,1]\to\eta$, 


$f:\{\tau(M\times E^n)\oplus(M\times E^{n+q})\}\times[0,1]\to\eta^{\perp}$ 


(\cite{Hm}, с.\,47). Ограничение $f$ на край 


$M\times\{0\}$ представляет собой изоморфизм 


$f_0:\tau(M\times E)\oplus(M\times E^{n+q})\to\nu_{2(n+q)}M$. 


Тогда 


$\xi=f_0\cdot(\left.\tau\right|_{M\times E^n}\oplus I_{M\times E^{n+q}})$, 


а оснащение пленки $M\times[0,1]$ задается с помощью $f$ и $g$. 


\end{pf*} 


 


Теорема 3 теперь вытекает непосредственно из 


теорем 1 и 2 и леммы 4. $\quad\square$ 


 


 


\begin{thebibliography}{9} 


 


\bibitem{Elw} 


Elworthy K.D., Tromba A.J. {\em Differential structures and Fredholm 


maps on Banach manifolds} // Proc. Symp. Pure Math. -- 1970. -- V.\,15. -- 


P.\,45--94. 


 


\bibitem{Zv1} 


Звягин В.Г., Ратинер Н.М. {\em Степень вполне непрерывных возмущений 


фредгольмовых отображений и ее приложение к бифуркации решений} // 


ДАН УССР. Сер.\,А. -- 1989. -- \No\,6. -- С.\,8--11. 


 


\bibitem{Zv2} 


Zvyagin V.G., Ratiner N.M. Oriented degree of Fredholm maps of 


non-negative index and its application to global bifurcation of 


solutions // Lect. Notes Math. -- 1992. -- \No 1520. -- P.111--137. 


 


\bibitem{Pon} 


Понтрягин Л.С. {\em Гладкие многообразия и их применения в теории 


гомотопий}. -- 2-е изд. -- М.: Наука, 1976. -- 175 с. 


 


\bibitem{Lin} 


Lindenstrauss J., Tzafriri L. {\em Classical Banach spaces}. I. 


{\em Sequence 


spaces} // Ergeb. Math. -- 1977. -- V.\,92. -- \No\,XIII. -- 190 р. 


 


\bibitem{Hm} 


Хьюзмоллер Д. {\em Расслоенные пространства}. -- М.: Мир, 1970. -- 442 с. 


 


\bibitem{Kra} 


Красносельский М.А. Забрейко П.П. {\em Геометрические методы нелинейного 


анализа}. -- М.: Наука, 1975. -- 511 с. 


 


\end{thebibliography} 


 


\vskip 2cm 


 


\post{\it Воронежский государственный университет}{\it Поступила} 


\post{}{30.01.1995} 


 


\end{document} 


